\documentclass[11pt]{article}
\usepackage[T1]{fontenc}
\usepackage{amsmath,amssymb}
\usepackage[margin=1in]{geometry}
\usepackage[hidelinks]{hyperref}
\hypersetup{pdftitle={Development summary: spherical Regge scissors congruence},
 pdfauthor={},pdfsubject={Retrospective account based on released project records}}
\ifdefined\pdftrailerid\pdftrailerid{}\fi
\setlength{\emergencystretch}{2em}
\title{Development summary:\\spherical Regge scissors congruence}
\author{}
\date{27 October 2026}
\begin{document}
\maketitle
\noindent\textbf{Discussions and model prompting: Alexander Kolpakov (UATX).}
This retrospective account describes the focused conditional manuscript and
its mathematical documentation. It is not a transcript or a reconstruction
of model turns. No original conversation export or authenticated model
identifier was supplied with the focused release.

\section*{Schematic research brief}
The target is to show that a nondegenerate spherical tetrahedron and its
elementary Regge mate admit finite geodesic dissections with pieces paired
by spherical isometries. A schematic brief is to identify the relevant
scissors-class defect, separate the finite algebra from the deeper
comparison theorems, and make every required input explicit. This is an
editorial description of the task, not a quotation of an original prompt.

The released result is conditional: the conclusion follows from the
numbered package E1--E8. The relative motivic and comparison bridge E2--E5
has not been constructed. The unconditional spherical Regge scissors
theorem is therefore not established by this contribution.

\section*{Mathematical route}
\begin{enumerate}
\item \textbf{Encode the elementary Regge operation.}
  The four moving lengths and angles transform by
  $R=\tfrac12J-I$, where $J$ is the all-ones matrix. Prove the matrix
  involution and orthogonality. Phase variables turn the corresponding
  angle transformation into $a_i\mapsto\tau/a_i$ with
  $\tau^2=\prod_i a_i$.
\item \textbf{Check the orientation-cover algebra.}
  Expand the Gram determinant, prove its relevant coefficient
  polynomials invariant under the phase substitution, and verify the
  two rational chart inverses and their cover on the stated open locus.
  Check the fixed-phase polynomial and its deck transformation.
\item \textbf{Expose the relative mathematical input.}
  Specify relative framed quadric coefficients, a strict three-weight
  setting, and the primitive-kernel comparison with
  $H^1(U,\mathbb Q(2))$. The proposed vanishing argument depends on
  the exact relative package; descriptions of its intended meaning do
  not supply that construction.
\item \textbf{Derive the finite channel cancellation.}
  The four moving channels cancel through $H^tH=4I$; the two fixed
  channels agree. Given the six-channel coproduct formula and descent
  in E4, the universal defect is primitive. E3 then forces it to vanish.
\item \textbf{Follow specialization to the spherical class.}
  Use the supplied commuting square, endpoint comparison, field-level
  injections and presentation comparison to kill the reduced scissors
  defect. Suspension, equal volume, area classification and cancellation
  then give a finite equidecomposition under the remaining inputs.
\item \textbf{Check the implication and its boundary.}
  Match every labeled manuscript item and every external-contract field
  to the formal declarations. Keep the unconstructed relative bridge
  explicit in the theorem statement and in the public account.
\end{enumerate}

\section*{Formal certificate and remaining work}
The Lean project checks the matrix, phase, determinant, chart, deck and
tensor algebra, followed by the kernel and scissors-group deductions. Its
main theorem takes a supplied \texttt{ProofContract} and returns a typed
finite-list equidecomposition certificate. Its abstract geometric
predicates and carriers are interpretation obligations; they do not build
actual spherical tetrahedra, motives or scissors groups from foundations.

The unresolved E2--E5 work includes constructing the relative face objects,
establishing the strict three-weight/primitive comparison, proving the
relative coproduct and Artin descent, and matching the specialized endpoint
with the field-level coefficient with the required signs and frames.
Published field-level results do not by themselves supply this exact
relative package. The Lean axiom audit checks the conditional implication,
not the existence of a geometric instance of its hypotheses.

The contributor supplied the discussion/model-prompting credit above.
This account authenticates no original dialogue or model run. Private
account and session metadata, credentials, local paths, detailed repair
logs and operational information are omitted.
\end{document}
